\section{Possessing the Truth Actively}

\paragraph{Decline and Cosmic Cycles}
In a recent comment, Jason-Adam brought up a few pertinent points that are too complex to address in a combox. For example, he wrote:

\begin{quotex}
One error Guenon and Evola might have made about the West I think is to judge it as a culture in formation against those who were fully fulfilled….looking into the future

I am confident in the future of the West because history shows all cultures have a Second Religiousness phase, where the springtime beliefs reemerge and overcome materialism and rationalism. 

\end{quotex}
I don't know that Guenon and Evola expressed it in quite those terms, but even so, it would be better to see the West as a ``culture in decline". That is, through a process of degeneration, it abandoned its roots in Tradition and has been implementing its opposite; specifically, it has rejected the sacred world while embracing the profane world. \textbf{Mircea Eliade} put it this way:

\begin{quotex}
all the definitions given up till now of the religious phenomenon have one thing in common: each has its own way of showing that the sacred and the religious life are the opposite of the profane and secular life 

\end{quotex}
So the ``second religiousness" that Jason-Adam mentioned can only be the return of the sacred and religious life. Once again, Eliade calls this the ``Eternal Return", that is, the mythic age returns, even as circumstances may differ. In terms of cosmic cycles, the end of the Kali Yuga, moving along the turn of the wheel, is followed by an era of increasing enlightenment. We see the Traditional societies of China and India increasingly following a mercantilist path, something Guenon had not anticipated. Chances are they will continue to decline into their own version of modernity and it will be the West that leads the return to Tradition.

\paragraph{The Age of the Holy Spirit}
\textbf{Blessed Joachim of Fiore}, whom Dante numbered among the elect, foresaw the age of the Holy Spirit following the Age of the Son. While care must be taken not to take this too literally, if understood as a deepening of understanding it makes sense. For the sacred and religious consciousness, it is the myth of the origin that dominates, not the end point. The Age of the Son revealed that the Logos is the origin of all. The consequences of this still need to be fully absorbed. The Franks understood this, a topic we will return to soon. 

Since the Owl of Minerva flies at dusk, we, at the decline of the West, are in a position to understand more directly, and more fully, what has happened and what is to be done. Those societies still at their peak cannot understand their possible future decline in an existential way. Unfortunately, we in the West are able to understand the modern world existentially, but seem to have forgotten any previous stage. That is why we take such pains in describing the most recent Traditional stage of the West: some few may be awakened to remember, and those few may turn into many.

\paragraph{Active Truth}
The Traditional teaching of the West, which may still be repeated today, claimed that the purpose of man in this life is to ``\emph{know, love, and serve God}." We can immediately see this as the basis of jnana, bhakti, and karma yogas for the West. To focus on jnana yoga for the moment, we need to point out that this goes well beyond the memorization and recitation of a catechism. The goal of the jnani is gnosis, or intellectual intuition, a direct knowing of God, beyond rational and literal formulations. In ``Modes of Spiritual Realization", \textbf{Frithjof Schuon} explains it this way:

\begin{quotex}
Only he who possesses the truth in an \textbf{active manner is really intellectual}, and not he who accepts it passively. 

\end{quotex}
He goes on to distinguish these two types, each of whom has ``learned a metaphysical truth":

\subparagraph{Active Type}
\begin{itemize}
\item He recognizes himself in it in a certain fashion 
\item He is capable of formulating it spontaneously 
\item He formulates in in an original and inspired manner 
\item He projects the light of his knowledge onto the most diverse contingencies 
\item This is because of a direct vision of the realities concerned and not by means of reasoning 
\end{itemize}
\subparagraph{Passive Type}
\begin{itemize}
\item He has a presentiment of its evidence 
\item He, however, is incapable of expressing it other than by repeating the doctrinal statement that was given to him 
\end{itemize}
\paragraph{Conclusion}
It is clearly important, or better said, it is the one thing necessary: that is, to know metaphysical truth in an active manner. This should require much contemplation of metaphysical teachings, and beyond that, a great deal of self-knowledge that is the fruit of self-observation. As long as a man remains passive in respect to metaphysical truths, he can only understand them as abstractions. This is revealed by his tendency to want to debate or choose sides in a dispute. Unable to directly intuit the ideas of Tradition, he thinks it has something to do with the past. Hence, he may try to ``prove" and ``demonstrate" these truths, for example, by false and tendentious historical studies. His efforts would be better spent in squaring the circle or inventing a perpetual motion machine. Such types are more dangerous than the totally ignorant.

We know many metaphysical truths. These must be pondered, they must be related to contingent conditions, whether in the past or the present. When a critical mass of men is achieved, then change will come. However, this is not a passive process, there is no historical necessity, but it depends instead on the collective and harmonious will of those who know.



\flrightit{Posted on 2012-11-26 by Cologero }

\begin{center}* * *\end{center}

\begin{footnotesize}\begin{sffamily}



\texttt{David on 2012-11-26 at 22:52 said: }

In vipassana meditation, this kind of active knowing is the highest form of wisdom (in the tradition of Goenka). When you speak of passive knowledge, it makes me think of submission to reason : you cling to your ideas, which is erroneous, since they are not "yours" to say. Debating or attaching yourself to those ideas is sign of passivity indeed. Very interesting read, as always.


\hfill

\texttt{Logres on 2012-11-27 at 13:23 said: }

Romanides was a tempting turn of thought, for me. I am glad I read Ruskin's praise of the Gothic, which seems to have been a genuine impulse of Soul in the vein of Tradition to resurrect (but not embalm or deify) what was good in the Classical soul. It's paradoxical because historically it ``looks" like a starting over, but am becoming convinced that it was a recovery, which may be what all ``conservative revolutions" (or the opposite of Revolution) looks like?


\hfill

\texttt{Jason-Adam on 2012-11-27 at 22:48 said: }

This is a very good article and I am pleased to have been the catalyst for Cologero to write this.

I am glad he mentioned the fact, which Evola had said too, that as the West is where the decline began, the West will also be the first to arise anew…..

Possessing the truth actively is something that is needed but is that possible without initiation ? In other words, can a person teach himself or does he need help from a master to enlighten himself ?


\hfill

\texttt{Jason-Adam on 2012-11-27 at 22:54 said: }

My remaining question of the period of the formation of the Franko-Roman-Gothic culture is, given the fact that it was a traditional society, what was the means by which Tradition was carried over from Ancient Rome to Mediaeval Rome ? Was it via Catholicism as Guenon thought or the Germanic Aryan blood as Evola thought or something else entirely ?


\hfill

\texttt{Caleb C on 2012-11-30 at 11:30 said: }

Both the priestly Catholic strain and the Germanic warrior strain came together powerfully when Charlemagne was crowned `Emperor of the Romans' by Pope Leo III on Christmas of 800. This was understood as a resurrection of the Western Roman Empire and transferring its imperium to the Frankish King. A beautiful example of how Tradition gives rebirth to the good and power of previous ages. The Holy Roman Empire lasted a good thousand years at the heart of Europe, the physical bodily manifestation of the Christian millennium until Napoleon, acting as an anti-christ of the counter-tradtion, killed it off in 1806.


\hfill

\texttt{Jason-Adam on 2012-11-30 at 14:30 said: }

Caleb C, your comment is absolutely stunning and overwhelming in its truth. If only more so called ``Traditionalists" would see instead of hearing…..


\end{sffamily}\end{footnotesize}
